\documentclass{beamer}
\usetheme[sidebar=false]{uu}
\usepackage{amsmath}
\usepackage{amssymb}
\usepackage{graphicx}

\usepackage{pgf}
\usepackage{tikz}
\usetikzlibrary{arrows, automata}
\usepackage[latin1]{inputenc}

\title[] {Functioneel Programmeren}
\subtitle{in vergelijking met andere paradigmas}
\author {Thomas van Binsbergen}
\date {}

\begin{document}
% new commands
\definecolor{brickred}{HTML}{b6321c}
\newcommand{\mvar}[1]{$\;\texttt{\textcolor{black}{#1}}\;$}
\newcommand{\mkey}[1]{$\;\texttt{\textcolor{orange}{#1}}\;$}
\newcommand{\mmod}[1]{$\;\texttt{\textcolor{purple}{#1}}\;$}
\newcommand{\mdef}[1]{$\;\texttt{\textcolor{blue}{#1}}\;$}
\newcommand{\mcon}[1]{$\;\texttt{\textcolor{brickred}{#1}}\;$}
\newcommand{\mtype}[1]{$\;\texttt{\textcolor{brown}{#1}}\;$}
\newcommand{\mpar}[1]{$(#1)$}
\newcommand{\mnil}[1]{\mcon{[]}}
\newcommand{\cons}[3]{\mpar{\mval{#1}\mcon{#2}\mval{#3}}}
\newcommand{\ra}[0]{$\;\rightarrow\;$}
\newcommand{\Ra}[0]{$\;\Rightarrow\;$}
\newcommand{\ind}[0]{$\;\;\;\;$}

\begin{frame}
    \titlepage
\end{frame}

\begin{frame}
    \frametitle{De paradigma's}
    \begin{itemize}
        \item Imperatief Programmeren (PHP, Javascript, Java, Python, Scala ...)
        \item Object Oriented programmeren (Java, PHP, Scala)
        \item Declaratief/Functioneel Programmeren (Haskell, Scala, ML)
    \end{itemize}
\end{frame}

\begin{frame}
    \frametitle{Programmeren zonder assignment}
    \begin{itemize}
        \item In de wiskunde zijn er veel variaties op het $=$-teken
        \item PHP en Javascript zouden het $:=$ symbool moeten gebruiken
        \item De zogenoemde 'destructieve assignment'
        \item In Haskell bestaat deze operator niet
        \item Daardoor hebben we geen variabelen!
    \end{itemize}
    \includegraphics[width=0.2\textwidth]{no_assignment.png}
\end{frame}

\begin{frame}
    \frametitle{Wat hebben we niet?}
    \begin{itemize}
        \item Variabelen
        \item Arrays (niets anders dan een variabele met structuur)
        \item Objecten (structuren)
        \item For en while loops (om door structuren te reizen)
    \end{itemize}
\end{frame}

\begin{frame}
    \frametitle{Wat hebben we wel?}
    \begin{itemize}
        \item Functies met parameters
        \item Functies zonder parameters, constanten
        \item Datatypes (structuren)
        \item Recursie (om door structuren te reizen)
    \end{itemize}
    \includegraphics[width=0.2\textwidth]{no_assignment.png}
\end{frame}

\end{document}
